\section[B]{Practicalities}\label{installation}
\subsection*{System Requirements}
 
You will need to have the following:

\begin{itemize}
\item{A computer.}
\item{A \GTAP 9 license and data file.}
\item{A \GAMS system, Version 24.7.1 or newer.}
\item{A nonlinear solver (PATH, IPOPT or CONOPT).}
\item{The MPSGE subsystem (optional)}
\end{itemize}
 
\subsection*{Getting Started}
 
The \GTAPINGAMS package is distributed as a zip file containing the
directory structure and \GAMS programs which can be unzipped into a
clean working directory. \GTAP source data are not distributed with
the \GTAPINGAMS system.  In order to generate models with the
\GTAPINGAMS tools, it is necessary to obtain the \GTAP 9 distribution
archives.
 
Here are the steps involved in installing \GTAPINGAMS:
\begin{enumerate}
\item   Create an empty root directory for \GTAP.
\item   Unzip \texttt{gtap9ingams.zip} in this directory.
\item   Install the \GTAP data file \texttt{flexagg9aY11.zip} into the \texttt{gtapdata}
subdirectory.\footnote{You need only one data files, but these data
are provided for three different years in the GTAP ``flexhar'' distribution.}
\item   Connect to the \texttt{build} directory and execute
\texttt{gams build.gms}  This will take some time to complete.  The batch
file is set up to begin with the 2011 base year data. Choose a different year if you wish.
\item   Run \texttt{test.gms} in the \texttt{build} directory  to 
evaluate benchmark consistency of the generated datasets.
\item   The \GAMS programs in the \texttt{forensics} directory provide
a template for how to conduct counterfactual policy simulations, using
Excel PivotTables and PivotCharts to synthesize model output.
\end{enumerate}
